\chapter{水中的氧}\label{ch:g7:oxygen-in-the-water}

山间溪流与夏日池塘都是“水”——
然而在其中一处游动的鳟鱼，到另一处就会喘不过气，
而在池塘里生长良好的鲤鱼，若来到溪流中，
则会发现溪流中找不到它的食物。两种水体的差异
很大程度上取决于一个看不见的数字：水中溶有多少
\omterm{def:g7:what-is-respiration:respiration}{氧气}。本章将学习读取、测量和预测这一
数字，以及由此描绘出的水生\omterm{def:g6:exploring-our-environment:components}{生物}分布图。

\section{溶解氧：稀缺的储备}

\begin{definition}[溶解氧]\label{def:g7:oxygen-in-the-water:dissolved}
水容纳氧的方式，就像甜茶容纳糖一样——
无形地\emph{溶解}\index{溶解氧}于其中。这一
数量极少：一升充分曝气的水中约含
\qty{10}{mg}\omterm{def:g7:what-is-respiration:respiration}{氧气}，比一升空气少约三十倍。
所有鳃\omterm{def:g7:what-is-respiration:respiration}{呼吸}及水下\omterm{def:g7:what-is-respiration:respiration}{呼吸}的\omterm{def:g6:exploring-our-environment:components}{生物}（见
\cref{ch:g7:how-animals-breathe}）都依赖于这一稀薄的
储备（见\cref{prop:g7:what-is-respiration:habitats}）。
\end{definition}

\begin{proposition}[补充溶解氧的来源]\label{prop:g7:oxygen-in-the-water:fills}
补充水中\omterm{def:g7:what-is-respiration:respiration}{氧气}的来源有两种：
\begin{enumerate}
  \item \emph{上方的空气}：\omterm{def:g7:what-is-respiration:respiration}{氧气}在水面溶入，
        在搅动、拍打或起泡之处溶入更快
        （急流、瀑布、波浪、
        雨）；
  \item \emph{下方的\omterm{def:g5:food-webs:roles}{生产者}}：白天，水生\omterm{def:g1:plants-around-us:plant}{植物}
        和绿色\omterm{def:g6:cell-unit-of-life:unicellular}{单细胞}\omterm{def:g6:exploring-our-environment:components}{生物}
        在制造\omterm{def:g6:plants-are-producers:matter}{有机物}时向水中释放\omterm{def:g7:what-is-respiration:respiration}{氧气}
        （见\cref{rem:g6:plants-are-producers:gas}）。
\end{enumerate}
\end{proposition}
\admitted

\begin{proposition}[氧的消耗与水的容纳能力]\label{prop:g7:oxygen-in-the-water:drains}
氧储备被\emph{每一种\omterm{def:g7:what-is-respiration:respiration}{呼吸}者}消耗：\omterm{def:g1:animals-around-us:animal}{动物}、
夜间的\omterm{def:g1:plants-around-us:plant}{植物}，尤其是\omterm{def:g6:decomposers-and-soil:decomposers}{分解者}，当有大量
死亡物质可供其取食时的
（见\cref{pb:g7:what-is-respiration:1}中的第八个瓶子）。上限也会变化：
\emph{冷}水能容纳的\omterm{def:g7:oxygen-in-the-water:dissolved}{溶解氧}明显多于温水——
例如，一条处于
\qty{5}{\degreeCelsius}的溪流所含\omterm{def:g7:oxygen-in-the-water:dissolved}{溶解氧}，
比\qty{25}{\degreeCelsius}的池塘多一半。\omterm{def:g6:exploring-our-environment:conditions}{温度}升高带来双重
挤压：能容纳的氧减少，而每一次
\omterm{def:g7:what-is-respiration:respiration}{呼吸}活动（见\cref{ex:g7:what-is-respiration:demand}）都进行得更快。
\end{proposition}
\admitted

\begin{example}[溪流与池塘对比]\label{ex:g7:oxygen-in-the-water:streampond}
现在，两种水体读起来就像两本账目。山间溪流：
水温低（上限高），每块卵石上都翻涌出白色泡沫，
（补充快），全年\omterm{def:g7:what-is-respiration:respiration}{氧气}都接近上限。夏季池塘：
水温高（上限低），水流静止（补充慢），
富含淤泥和死亡物质（\omterm{def:g6:decomposers-and-soil:decomposers}{分解者}食量大）——
\omterm{def:g7:what-is-respiration:respiration}{氧气}几乎贴着池底，最差的时候是凌晨，
黎明前，此时\omterm{def:g1:plants-around-us:plant}{植物}已整夜进行\omterm{def:g7:what-is-respiration:respiration}{呼吸作用}，
而没有白天的制造来加以抵消。
\end{example}

\begin{omfigure}
\includegraphics[width=0.62\linewidth]{images/book1/ai/g7-oxygen-stream.jpg}

{\small 富氧的一端：冷水，在每块卵石处翻涌成白色——
水温高、上限高，且持续补充。}
\end{omfigure}

\begin{omfigure}
\begin{tikzpicture}
\begin{axis}[omaxis,
    width=9.6cm, height=5.6cm,
    xlabel={一天中的时刻}, ylabel={溶解氧 (mg/L)},
    xmin=0, xmax=27, ymin=0, ymax=14,
    xtick={0,6,12,18,24},
    ytick={4,8,12},
    ]
  % stream: flat high
  \addplot[very thick, omDef] coordinates {
    (0,10.6) (6,10.7) (12,10.4) (18,10.5) (24,10.6)
  };
  % pond: sinusoidal-ish, low before dawn, peak late afternoon
  \addplot[very thick, omThm, smooth] coordinates {
    (0,4.4) (3,3.4) (6,3.0) (9,5.2) (12,7.8) (15,9.4)
    (17,9.8) (20,7.4) (24,4.6)
  };
  \node[font=\footnotesize, omDef] at (axis cs:4.6,12.1)
    {山间溪流};
  \node[font=\footnotesize, omThm] at (axis cs:8.2,1.6)
    {夏季池塘};
\end{axis}
\end{tikzpicture}

{\small 一天中的\omterm{def:g7:oxygen-in-the-water:dissolved}{溶解氧}变化。冷而翻涌的溪流
保持在较高上限附近稳定；温暖的池塘
依赖其\omterm{def:g5:food-webs:roles}{生产者}——白天不断上升，
夜间下降至黎明时的最低点。}
\end{omfigure}

\begin{example}[读懂池塘的曲线]\label{ex:g7:oxygen-in-the-water:curve}
池塘的日变化，是
\cref{rem:g6:plants-are-producers:gas}用数据呈现的样子：日出
启动了\omterm{def:g5:food-webs:roles}{生产者}的制造，曲线随之上升；
傍晚时达到峰值；夜间，每一个\omterm{def:g7:what-is-respiration:respiration}{呼吸}者
都消耗这一储备，唯一的补充来源只剩静止的水面
——曲线滑落至黎明时的最低点。\omterm{prop:g3:sorting-living-things:groups}{鱼类}
死亡在过度繁茂的池塘中常发生于黎明，现在你知道原因了。
\end{example}

\section{栖居生物的分布图}

\begin{proposition}[氧气绘制分布图]\label{prop:g7:oxygen-in-the-water:map}
由于各\omterm{def:g5:biodiversity:species}{物种}的需氧量各不相同，\omterm{def:g7:oxygen-in-the-water:dissolved}{溶解
氧}读数便勾勒出栖居\omterm{def:g6:exploring-our-environment:components}{生物}的分布
（见\cref{prop:g6:exploring-our-environment:decide}，
适用于水体）：
\begin{itemize}
  \item \emph{需氧苛刻的}\omterm{def:g5:biodiversity:species}{物种}——鳟鱼、蜉蝣幼虫
        带着摆动的鳃丝束
        （见\cref{pb:g7:how-animals-breathe:1}）——仅生活在
        寒冷、翻涌、富氧的水中；
  \item \emph{耐受性强的}\omterm{def:g5:biodiversity:species}{物种}——鲤鱼、池塘螺、
        鼠尾蛆幼虫以伸向水面的\omterm{def:g7:what-is-respiration:respiration}{呼吸}管——
        在温暖、静止、贫氧的水中也能生存；
  \item 因此普查须倒过来读：若在石块下找到蜉蝣幼虫，
        说明水质良好；若只找到
        \omterm{def:g7:what-is-respiration:respiration}{呼吸}管型和吞气型\omterm{def:g6:exploring-our-environment:components}{生物}，则水质较差。
\end{itemize}
河流管理员对小型\omterm{def:g6:exploring-our-environment:components}{生物}进行普查，正是出于这一
原因：栖居\omterm{def:g6:exploring-our-environment:components}{生物}记录的是水体近期最糟的几天，
而一次探针读数可能恰好漏掉这些时段。
\end{proposition}
\admitted

\begin{example}[有机废水的故事]\label{ex:g7:oxygen-in-the-water:waste}
某村庄未经处理的废水——富含\omterm{def:g6:plants-are-producers:matter}{有机物}——
流入一条清澈的小溪。排污管下游，
\omterm{def:g6:decomposers-and-soil:decomposers}{分解者}大量摄食并繁殖
（见\cref{def:g6:decomposers-and-soil:decomposers}），它们
大量\omterm{def:g7:what-is-respiration:respiration}{呼吸作用}耗尽了\omterm{def:g7:what-is-respiration:respiration}{氧气}：鳟鱼和蜉蝣
幼虫消失，耐受性强的\omterm{def:g7:what-is-respiration:respiration}{呼吸}管型\omterm{def:g6:exploring-our-environment:components}{生物}取而代之，
再往下游——废物已被分解，水体
经浅滩复氧——需氧苛刻的\omterm{def:g5:biodiversity:species}{物种}重新出现。
污染物是\omterm{def:g6:plants-are-producers:matter}{有机物}，真正的杀手是
缺氧；分布图把这一切都记录了下来。
\end{example}

\begin{method}[判断水体]\label{met:g7:oxygen-in-the-water:judge}
判断溪流或池塘的水质，可按以下步骤：
\begin{enumerate}
  \item 测量或估算供氧因素：水温、
        水体翻动程度、绿色\omterm{def:g1:plants-around-us:plant}{植物}的生长；
  \item 评估耗氧因素：淤泥与死亡\omterm{def:g6:plants-are-producers:matter}{有机物}、
        \omterm{def:g6:exploring-our-environment:components}{生物}拥挤程度、任何有机废水的输入；
  \item 普查石块下及水草间的小型\omterm{def:g6:exploring-our-environment:components}{生物}
——是否有需氧苛刻的\omterm{def:g5:biodiversity:species}{物种}，还是仅有耐受性强的\omterm{def:g5:biodiversity:species}{物种}？
  \item 若在某一时刻取样，须记住池塘的曲线：
        晴好午后的读数可能掩盖致命的黎明。
\end{enumerate}
\end{method}

\section{练习}

\begin{exercise}[$\star$]\label{exo:g7:oxygen-in-the-water:1}
估算一升充分曝气的水中大约含有多少\omterm{def:g7:what-is-respiration:respiration}{氧气}，
并与空气相比如何？
\end{exercise}

\begin{exercise}[$\star$]\label{exo:g7:oxygen-in-the-water:2}
写出水体中\omterm{def:g7:what-is-respiration:respiration}{氧气}的两个来源，以及
各来源在何种\omterm{def:g6:exploring-our-environment:conditions}{条件}下供氧最快。
\end{exercise}

\begin{exercise}[$\star$]\label{exo:g7:oxygen-in-the-water:3}
写出氧储备的三个消耗因素。哪一个
在死亡\omterm{def:g6:plants-are-producers:matter}{有机物}大量存在时会急剧增大？
\end{exercise}

\begin{exercise}[$\star$]\label{exo:g7:oxygen-in-the-water:4}
请说明\omterm{def:g6:exploring-our-environment:conditions}{温度}对水中\omterm{def:g7:what-is-respiration:respiration}{呼吸}\omterm{def:g6:exploring-our-environment:components}{生物}施加的双重压力。
\end{exercise}

\begin{exercise}[$\star$]\label{exo:g7:oxygen-in-the-water:5}
为什么池塘的氧含量曲线在傍晚时达到峰值，
并在黎明时降至最低？
\end{exercise}

\begin{exercise}[$\star$]\label{exo:g7:oxygen-in-the-water:6}
哪两种普查结果标志着富氧水体，
哪两种标志着水质较差？
\end{exercise}

\begin{exercise}[$\star\star$]\label{exo:g7:oxygen-in-the-water:7}
请向池塘主人解释
\cref{ex:g7:oxygen-in-the-water:curve}中黎明时的\omterm{prop:g3:sorting-living-things:groups}{鱼类}死亡——
并从每一个供氧来源出发，各提出一项补救措施，
依据\cref{prop:g7:oxygen-in-the-water:fills}。
\end{exercise}

\begin{exercise}[$\star\star$]\label{exo:g7:oxygen-in-the-water:8}
将\cref{ex:g7:oxygen-in-the-water:waste}以链条形式复述：
废物输入，随后发生什么，再随后发生什么——
以及溪流为何在下游得以恢复？
\end{exercise}

\begin{exercise}[$\star\star$]\label{exo:g7:oxygen-in-the-water:9}
为什么河流管理员更信赖小型\omterm{def:g6:exploring-our-environment:components}{生物}的普查，
而非单次探针读数？请运用“近期最糟的几天”这一思想
作答。
\end{exercise}

\begin{exercise}[$\star\star$]\label{exo:g7:oxygen-in-the-water:10}
庭院池塘变成豌豆汤般的绿色，随后经历
一周无风、炎热的天气。请借助两条命题追踪危险——
并说明池塘主人应在何时将喷泉开到最大。
\end{exercise}

\begin{exercise}[$\star\star$]\label{exo:g7:oxygen-in-the-water:11}
鳟鱼养殖场建在寒冷河流上，配有堰坝与
溅水式养殖渠，从不建在温暖、静止的湖泊上。请
逐一论证这句话中的每个分句。
\end{exercise}

\begin{exercise}[$\star\star\star$]\label{exo:g7:oxygen-in-the-water:12}
“水中，\omterm{def:g7:what-is-respiration:respiration}{氧气}是货币；\omterm{def:g6:exploring-our-environment:conditions}{温度}决定货币供应量；
\omterm{def:g6:decomposers-and-soil:decomposers}{分解者}则是最大的开销大户。”
请将每个分句译成本章的通俗论断，
并为每句各引用一条命题或一个例子。
\end{exercise}

\section{问题：河道看守人的年鉴}

\begin{problem}[{周末问题——一条河、四个站点、
一年的读数}]\label{pb:g7:oxygen-in-the-water:1}
一名河道看守人为四个站点各存一本年鉴：K1，
堰坝下游的寒冷急流；K2，流经牧场的
缓慢温暖河段；K3，紧邻乳品厂的废水管下方；K4，再往下游两公里，
越过一段多石的长浅滩。每个站点
都要记录探头读数，并进行一次小型\omterm{def:g6:exploring-our-environment:components}{生物}普查。

\medskip\noindent\textbf{第一部分——解读各站点数据。}
\begin{enumerate}
  \item K1 读数为 \qty{11}{mg/L}，
        水温 \qty{8}{\degreeCelsius}；K2 读数为 \qty{7}{mg/L}，
        水温 \qty{22}{\degreeCelsius}。请将这一差异
        归因于恰当的\omterm{def:g7:what-is-respiration:respiration}{氧气}来源与上限。
  \item K1 的普查：鳟鱼苗、蜉蝣幼虫；K2 的：鲤鱼、
        池塘螺、鼠尾幼虫。请用\cref{prop:g7:oxygen-in-the-water:map}解释这两份普查结果。
        依据\cref{prop:g7:oxygen-in-the-water:map}。
  \item K3 读数为 \qty{3}{mg/L}，尽管水温
        不比 K2 高。请指出那处排水口，以及
        造成排放的工人。
  \item K4 读数为 \qty{9}{mg/L}，其普查结果再次显示
        蜉蝣幼虫。请从\omterm{def:g7:what-is-respiration:respiration}{氧气}来源和
        废物去向两方面解释其恢复。
\end{enumerate}

\medskip\noindent\textbf{第二部分——全年的节律。}
\begin{enumerate}[resume]
  \item 八月份 K2 的黎明读数降至
        \qty{4}{mg/L}，而其午后读数可达
        \qty{9}{mg/L}。请画出或描述其日变化曲线，
        并指出此消彼长的两个过程。
  \item K2 一月份的读数比八月更稳定、更高，
        尽管牧场河段依旧
        静止如故。请给出冬季的两个原因。
  \item 预报显示将有一段炎热无风的两周。请
        按\omterm{prop:g3:sorting-living-things:groups}{鱼类}死亡风险对四个站点排序，每个给出
        一条理由。
  \item 看守人将八月的探头测量定在
        恰好日出之前。请为这一安排辩护。
\end{enumerate}

\medskip\noindent\textbf{第三部分——结论与
补救措施。}
\begin{enumerate}[resume]
  \item 乳品厂安装了一套能去除大部分
        排放水中\omterm{def:g6:plants-are-producers:matter}{有机物}的处理设施。预测一年后 K3 的
        探头读数与普查结果，按恢复先后
        排列。
  \item 一位土地所有者提议砍伐 K2 的岸边树木，
        理由是“让更多光照产生更多\omterm{def:g7:what-is-respiration:respiration}{氧气}”。请权衡
        这一提议：遮荫还起什么作用，据
        \cref{prop:g7:oxygen-in-the-water:drains}，并
        说明你会如何回应。
  \item 普查在一个站点发现了蜉蝣幼虫，而该站点
        的单次午后探头读数看上去平平。
        你更相信哪一条证据能说明整个季节的
        情况？为什么？
  \item 以看守人为整条河流拟定的一句话规则结束年鉴，
        并使用\emph{冷}、
        \emph{搅动}和\emph{未受负荷}
        （指有机废物）。
\end{enumerate}
\end{problem}
